%=======================================================================
%  APPENDICES
%=======================================================================
\section*{Appendix A -- AI Use Statement}
\label{sec:appendixA}
\addcontentsline{toc}{section}{Appendix A -- AI Use Statement}
\setcounter{table}{0}
\renewcommand{\thetable}{A.\arabic{table}}

Table~\ref{tab:aiuse} lists the AI-supported activities in this project, what each
contributed and how we checked the output. No factual claim in the report is cited to
an AI tool; every biological and technical statement carries a reference of its own.

{\small
\renewcommand{\arraystretch}{1.15}
\begin{longtable}{L{2.4cm} L{3.2cm} L{4.0cm} L{4.4cm}}
\caption{AI use statement.}\label{tab:aiuse}\\
\toprule
\textbf{AI tool} & \textbf{Used for} & \textbf{How the output influenced the project} & \textbf{How important output was verified} \\
\midrule
\endfirsthead
\toprule
\textbf{AI tool} & \textbf{Used for} & \textbf{How the output influenced the project} & \textbf{How important output was verified} \\
\midrule
\endhead
\bottomrule
\endfoot

AskNature website and AskNature Chat &
Searching biological strategies for functions F1--F5 with the keywords and biologized
questions of Section~\ref{sec:researchpath} &
Produced the candidate list of Table~\ref{tab:models}, including the models we later
de-selected &
Each strategy page was opened and read. Where a primary paper exists it is cited in
Section~\ref{sec:models} instead of the AskNature page (silver ant, cactus, camel, snail,
pine cone, honeybee); the AskNature pages remain cited as the lead for the elephant models
and the Morpho \\

Large language model assistant (Claude, Anthropic) &
Structuring the scoping notes; drafting and editing report text; converting the report to
\LaTeX; an editorial audit of the draft against the course rubric (repetition, unsupported
claims, citation gaps); recomputing the engineering estimates &
Draft wording in Sections~1--5 and the layout of most tables; the audit produced the list
of corrections applied in this version (figure numbering, missing references, KPI
definitions, removal of repeated passages) &
The team decided on every audit finding. All numbers in Sections~\ref{sec:check1},
\ref{sec:check2} and \ref{sec:v2check} were recomputed by hand from the stated inputs
before the text was accepted. Any claim the assistant proposed without a source was
either given a source we read ourselves or removed \\

AI as critical reviewer &
Reviewing Concept~A~V1 and B~V1 with the Engineering Canvas, the functions and the KPIs
(Section~\ref{sec:review}) &
Four review issues in Table~\ref{tab:review}; three were accepted and changed the design
(pipes into external panels, reduced connected area, KPI units), one was qualified &
Each issue was assessed by the team against the literature named in
Section~\ref{sec:review}; the qualified issue (ventilation) was not accepted because the
cited studies disagree \\

Image generation &
\todo{State whether any figure was AI-generated (for example the concept renders of
Figures~\ref{fig:conceptA} and \ref{fig:v2loop}); if so, name the tool} &
\todo{Illustration only; no figure is used as measurement or evidence} &
\todo{Each generated image was compared with the hand sketch it was made from} \\
\end{longtable}}

\section*{Appendix B -- Optional Supporting Material}
\label{sec:appendixB}
\addcontentsline{toc}{section}{Appendix B -- Optional Supporting Material}
\renewcommand{\thetable}{B.\arabic{table}}
\setcounter{table}{0}

\subsection*{B.1 Glossary}

\begin{description}\setlength{\itemsep}{2pt}
  \item[Albedo] Refers to an index from 0 -- 1 on how much light is reflected.
  \item[PV] Photovoltaic
  \item[Radiation] Shortwave radiation is the type of radiation emitted by the sun.
    Longwave radiation is thermal radiation which is emitted from a surface into the
    surrounding atmosphere.
  \item[UHI] Urban Heat Island
  \item[UV-Light] Ultraviolet Light
\end{description}

\needspace{14\baselineskip}
\subsection*{B.2 Decision log}

Table~\ref{tab:decisionlog} lists the decisions that shaped the concept, the alternative
each displaced, and where in the report the reason is given.

{\small
\renewcommand{\arraystretch}{1.15}
\begin{longtable}{L{2.6cm} L{3.1cm} L{3.1cm} L{5.2cm}}
\caption{Decision log.}\label{tab:decisionlog}\\
\toprule
\textbf{Decision} & \textbf{Chosen} & \textbf{Rejected alternative} & \textbf{Reason and where it is argued} \\
\midrule
\endfirsthead
\toprule
\textbf{Decision} & \textbf{Chosen} & \textbf{Rejected alternative} & \textbf{Reason and where it is argued} \\
\midrule
\endhead
\bottomrule
\endfoot
Design question & Reduce the heat island by changing façade and roof of existing buildings & Three broader or wrongly framed candidates & Anchored to the system boundary and testable; Section~\ref{sec:designquestion} \\
Reference city and site & Zurich, Modellierungsgebiet 03 & A generic city or a freely chosen building & Documented heat map and albedo classes, open geodata; Section~\ref{sec:context} \\
System of interest & Façade and roof of the building envelope & The whole building, internal cooling & Keeps the project finishable; Table~\ref{tab:inout} \\
KPI units & Physical units with baseline and target & Temperature labels without targets & Needed for a defensible comparison; Section~\ref{sec:kpi} \\
Reflective surface & On the roof & On the façade (Concept C) & Reflective walls raised pedestrian heat stress in canyon simulations; Section~\ref{sec:morph} \\
Water and power & No evaporating layer, no fans & Wet layer, fan-driven cavity & Water and Energy constraints of Table~\ref{tab:constraints}; Section~\ref{sec:morph} \\
Concept pair & A (water façade) and B (textured coating) & C (responsive louvers) & Humidity is not a cooling signal; Section~\ref{sec:morph} \\
Final concept & A, revised to V2 & B & B's benefit is unproven and of uncertain sign; Section~\ref{sec:review} \\
Pipe location & Separable external collector panels with a drainage gap & Absorber layer without access (V1) & Hidden damage and no repair access; Section~\ref{sec:v2} \\
Fins at ground level & Start at the first floor, thinner base & Fins to the ground (V1) & Pavement width and door operation; Section~\ref{sec:v2access} \\
Roof network & Many short parallel circuits & Four long coils & Long coils lose about 67 times the available buoyancy pressure; Section~\ref{sec:v2check} \\
Frost protection & Drainback & Glycol & Keeps an additive out of a loop on an occupied façade, at the cost of shoulder-season operation; Sections~\ref{sec:winter} and \ref{sec:benefits} \\
\end{longtable}}
